\chapter{Kết quả ứng dụng}
\label{chap:results}

\section{Đặc tả yêu cầu phần mềm}

\subsection{Mô tả tổng quan}

PERFIN là nguyên mẫu ứng dụng di động phục vụ nhập, chuẩn hóa, quản lý và phân tích dữ liệu tài chính cá nhân. Trọng tâm của nguyên mẫu là bảo vệ chất lượng dữ liệu và kiểm chứng các giải thuật; giao diện chat chỉ là một cổng nhập bổ sung cho các màn hình trực tiếp. Mọi phép tính tài chính, kiểm tra ràng buộc và thay đổi dữ liệu đều do backend thực hiện. LLM, OCR và STT không truy cập trực tiếp PostgreSQL và không được xem là nguồn số liệu chuẩn.

Phiên bản hiện tại vận hành với một hồ sơ \texttt{default\_user}. Các khóa ngoại theo người dùng tạo đường nâng cấp về sau nhưng chưa chứng minh hệ thống có xác thực, phân quyền hoặc cô lập nhiều người dùng ở mức production. PostgreSQL là nguồn dữ liệu nghiệp vụ bền vững; Redis lưu trạng thái có TTL, cache và dữ liệu queue, với bộ nhớ tiến trình làm fallback trong môi trường phát triển. Phần đặc tả dùng mã FR/NFR, điều kiện chấp nhận và mã ca kiểm thử để tạo khả năng truy vết theo tinh thần IEEE~830 \cite{ieee830}.

\begin{table}[H]
  \centering
  \caption{Tác nhân của hệ thống}
  \label{tab:actors}
  \begin{tabularx}{\textwidth}{@{}L{3.4cm}Y Y@{}}
    \toprule
    \textbf{Tác nhân} & \textbf{Vai trò} & \textbf{Giới hạn} \\
    \midrule
    Người dùng & Nhập dữ liệu, sửa/xác nhận/hủy bản xem trước, quản lý ngân sách/mục tiêu và xem báo cáo. & Chịu trách nhiệm kiểm tra dữ liệu suy ra từ ngôn ngữ hoặc media trước khi lưu. \\
    Ứng dụng di động & Thu nhận văn bản, ảnh, âm thanh; gọi REST API; hiển thị bản xem trước, dữ kiện và kết quả. & Không trực tiếp truy cập cơ sở dữ liệu hoặc giữ khóa bí mật của nhà cung cấp AI. \\
    API và dịch vụ nghiệp vụ & Validation, transaction cơ sở dữ liệu, tính toán và trả kết quả có cấu trúc. & Không giao phép tính hoặc quyền ghi dữ liệu cho LLM. \\
    Dịch vụ AI ngoài & Trả lời gọi công cụ, văn bản OCR hoặc transcript STT. & Kết quả là đầu vào chưa tin cậy; không tự thực thi nghiệp vụ. \\
    Worker nền & Chạy nhắc định kỳ, quét runway/subscription, insight cuối tháng và dọn export. & Phải giới hạn theo hồ sơ và chống hiệu ứng lặp khi retry. \\
    Quản trị viên phát triển & Cấu hình provider, migration và dữ liệu demo. & Không phải tác nhân người dùng của sản phẩm. \\
    \bottomrule
  \end{tabularx}
\end{table}

Các mã \texttt{TC-FRxx-yy} dưới đây là định danh ca kiểm thử ở mức đặc tả. Một ca chỉ được xem là ``đã đo'' khi có dữ liệu đầu vào, kết quả mong đợi, môi trường và log tương ứng tại mục 3.3.

\subsection{Yêu cầu cụ thể}

\subsubsection{Giao diện và phụ thuộc bên ngoài}

\begin{itemize}
  \item Ứng dụng di động giao tiếp với backend Express qua REST/JSON; ứng dụng không kết nối trực tiếp PostgreSQL, Redis hoặc API của nhà cung cấp AI.
  \item Đầu vào media chấp nhận các định dạng ảnh/audio do backend cho phép, tối đa 10~MB mỗi tệp. Transcript hoặc raw text phải được trả để người dùng kiểm tra.
  \item Gemini hoặc parser cục bộ đảm nhiệm hiểu ý định và trích xuất trường; PaddleOCR/Google Vision và PhoWhisper/Google Speech là các adapter có thể thay thế. Provider không khả dụng phải dẫn đến lỗi hoặc fallback thật, không dựng kết quả mẫu.
  \item PostgreSQL giữ dữ liệu nghiệp vụ và lịch sử cần kiểm toán. Redis giữ pending/clarification, cache và queue; mất Redis không được làm phát sinh dữ liệu tài chính giả.
\end{itemize}

\subsubsection{Quy tắc dữ liệu và nghiệp vụ dùng chung}

\begin{enumerate}
  \item Số tiền giao dịch thu--chi, chuyển ví, ngân sách và khoản định kỳ phải là số hữu hạn dương. Giá trị lãi/lỗ đầu tư là số hữu hạn khác 0 và có thể mang dấu.
  \item Ngày giao dịch, chuyển ví và thanh toán không được ở tương lai; ngày mục tiêu không được ở quá khứ khi tạo kế hoạch.
  \item Danh mục và ví tham chiếu phải tồn tại trong phạm vi hồ sơ hiện tại. Phiên bản demo chưa có authorization nhiều người dùng.
  \item Số dư ví \textbf{được phép âm}. PERFIN ghi sổ và cảnh báo dòng tiền, không từ chối chi/chuyển chỉ vì số dư nguồn không đủ.
  \item Thay đổi do chat hoặc media khởi tạo phải qua bản xem trước và xác nhận; thao tác trực tiếp từ form có thể ghi sau khi backend validation thành công.
  \item Thao tác ảnh hưởng nhiều bản ghi hoặc số dư phải nằm trong một transaction cơ sở dữ liệu. Lỗi cache/hydration sau \texttt{COMMIT} không được báo thành lỗi ghi khiến người dùng retry.
  \item Giao dịch dùng soft delete và có thể khôi phục; báo cáo mặc định không tính bản ghi đã xóa.
  \item Kết quả phân tích phải có kỳ dữ liệu, dữ kiện định lượng và phương pháp. Thiếu dữ liệu phải trả \texttt{null}/cảnh báo, không suy diễn kết luận mạnh.
\end{enumerate}

\subsubsection{Phạm vi và mức ưu tiên}

\begin{table}[H]
  \centering
  \caption{Phân nhóm chức năng theo mức ưu tiên học thuật}
  \label{tab:function-groups}
  \begin{tabularx}{\textwidth}{@{}L{3.4cm}Y Y@{}}
    \toprule
    \textbf{Nhóm} & \textbf{Chức năng} & \textbf{Vai trò trong niên luận} \\
    \midrule
    Lõi dữ liệu & Giao dịch, ví, danh mục, ngân sách, khoản định kỳ và mục tiêu. & Đối tượng nghiên cứu chính. \\
    Lõi giải thuật & Parsing, matching, feedback, analytics, budget/goal planner. & Đóng góp chính cần mô tả và kiểm thử. \\
    Lớp LLM & Tool routing, extraction, clarification, narration và persona. & Thành phần hỗ trợ có ranh giới. \\
    Hỗ trợ demo & Dashboard, màn hình quản lý và export. & Chứng minh khả năng sử dụng, không phải trọng tâm. \\
    Ngoài phạm vi & Auth production, bank sync, shared wallet và high availability. & Hướng phát triển, không phải điều kiện hoàn thành niên luận. \\
    \bottomrule
  \end{tabularx}
\end{table}

\subsection{Yêu cầu chức năng}

\begin{longtable}{@{}L{1.3cm}L{4.3cm}L{8.2cm}@{}}
  \caption{Yêu cầu chức năng}\label{tab:functional-requirements}\\
  \toprule
  \textbf{Mã} & \textbf{Yêu cầu} & \textbf{Điều kiện chấp nhận cấp cao} \\
  \midrule
  \endfirsthead
  \toprule
  \textbf{Mã} & \textbf{Yêu cầu} & \textbf{Điều kiện chấp nhận cấp cao} \\
  \midrule
  \endhead
  FR-01 & Nhập giao dịch bằng văn bản tự nhiên & Chuẩn hóa một hoặc nhiều giao dịch, tạo bản xem trước và không ghi khi thiếu trường bắt buộc. \\
  FR-02 & Nhập bằng ảnh và giọng nói & Trả raw OCR/transcript thật; xác nhận trước pipeline giao dịch; chưa tuyên bố đối soát tổng hóa đơn. \\
  FR-03 & Clarification và giao dịch chờ & Lưu state 5 phút; cho phép bổ sung/sửa/xác nhận/hủy; một preview chỉ được claim một lần. \\
  FR-04 & Phân loại và học từ phản hồi & Ưu tiên exact/alias/fuzzy an toàn; ghi correction sau commit; không học khi kết quả mơ hồ/xung đột. \\
  FR-05 & Quản lý dữ liệu tài chính & CRUD có validation và soft delete; cập nhật số dư nguyên tử; cho phép số dư ví âm. \\
  FR-06 & Chuyển ví nguyên tử & Debit, credit và lịch sử cùng transaction; hai ví khác nhau; không kiểm tra đủ số dư. \\
  FR-07 & Phân tích dữ liệu & Tính trend, anomaly, runway, recurring và tương quan dương bằng giải thuật xác định. \\
  FR-08 & Insight có căn cứ và persona & Trả facts và lời diễn giải; numeric grounding checker là thiết kế đích, chưa phải kết quả đã đo. \\
  FR-09 & Ngân sách và dự báo & Tính tiến độ/dự báo; đề xuất theo lịch sử hoặc khung tỷ lệ; chỉ áp dụng khi xác nhận. \\
  FR-10 & Mục tiêu tài chính & Preview kế hoạch tiết kiệm/mua sắm/trả nợ và what-if; chỉ lưu đúng payload đã xem trước. \\
  FR-11 & Khoản định kỳ và tác vụ chủ động & Quản lý lịch, thanh toán nguyên tử; worker retry có dedup và user scope. \\
  FR-12 & Xuất dữ liệu & CSV hoạt động; luồng mang nhãn PDF hiện tạo HTML; lưu lịch sử và dọn file hết TTL. \\
  \bottomrule
\end{longtable}

\subsubsection{FR-01 --- Nhập giao dịch bằng văn bản tự nhiên}

\textbf{Mục đích và kích hoạt.} Giảm nhập tay nhưng vẫn bảo đảm schema. Người dùng gửi một câu chat hoặc gọi endpoint phân tích; câu có thể chứa một hay nhiều giao dịch.

\textbf{Đầu vào và xử lý.} Văn bản không rỗng; hồ sơ có ví và danh mục. Trường đích gồm mô tả, amount dương, loại \texttt{income/expense}, ngày không ở tương lai, danh mục và ví. Orchestrator chọn LLM structured output hoặc parser cục bộ, chuẩn hóa tiền/ngày theo mục 2.2.1 và ánh xạ danh mục/ví. Backend áp dụng quy tắc sổ cái mục 2.1; LLM chỉ đề xuất đối số. Batch chỉ được ghi nguyên tử sau bản xem trước FR-03.

\textbf{Lỗi, đầu ra và hậu điều kiện.} Thiếu amount hoặc còn mơ hồ chuyển sang clarification; provider lỗi dùng fallback thật hoặc trả lỗi. Ngày tương lai, tham chiếu ngoài phạm vi hoặc một phần tử batch sai làm từ chối toàn bộ. Trước xác nhận chỉ trả dữ liệu chuẩn hóa, cảnh báo và \texttt{pending\_id}; sau xác nhận mới đổi PostgreSQL, tiêu thụ pending và vô hiệu cache best-effort.

\textbf{Tiêu chí chấp nhận.}
\begin{itemize}
  \item \texttt{TC-FR01-01}: câu đủ trường tạo một preview và chưa tăng số giao dịch trong DB.
  \item \texttt{TC-FR01-02}: câu nhiều giao dịch trích xuất đủ và commit tất-cả-hoặc-không.
  \item \texttt{TC-FR01-03}: thiếu amount tạo clarification, không tạo pending có thể ghi sai.
  \item \texttt{TC-FR01-04}: ngày tương lai hoặc tham chiếu ngoài phạm vi không đổi số dư.
\end{itemize}

\subsubsection{FR-02 --- Nhập giao dịch bằng ảnh và giọng nói}

\textbf{Mục đích và kích hoạt.} Chuyển ảnh hóa đơn hoặc âm thanh tiếng Việt thành văn bản có thể kiểm tra rồi tái sử dụng FR-01. Tệp phải thuộc loại backend hỗ trợ và không quá 10~MB.

\textbf{Đầu vào và xử lý.} Adapter media phải được cấu hình hoặc có provider cục bộ. Provider trả raw OCR/transcript kèm tên provider; hệ thống hiển thị kết quả, trích xuất trường rồi tạo lựa chọn/bản xem trước. Nền tảng OCR/STT ở mục 2.2.4; CER/WER và độ đúng trường ở mục 2.4.2. Phiên bản hiện tại có thể cho chọn lưu tổng hóa đơn hoặc từng dòng, nhưng \textbf{chưa đối soát tổng mặt hàng với tổng hóa đơn}.

\textbf{Lỗi, đầu ra và hậu điều kiện.} File sai loại/quá cỡ hoặc upload lỗi bị từ chối. Provider không trả văn bản phải sinh lỗi thật, không tạo mock. API trả raw text/transcript, provider, ngữ cảnh và dữ liệu trích xuất; chỉ sau xác nhận media và FR-03 mới ghi giao dịch.

\textbf{Tiêu chí chấp nhận.}
\begin{itemize}
  \item \texttt{TC-FR02-01}: ảnh hợp lệ trả raw OCR và provider, không tuyên bố tổng dòng hàng đã khớp.
  \item \texttt{TC-FR02-02}: transcript được hiển thị, sửa/xác nhận rồi mới vào parser.
  \item \texttt{TC-FR02-03}: provider lỗi không tạo mock hoặc pending.
  \item \texttt{TC-FR02-04}: file sai loại hoặc lớn hơn 10~MB bị chặn trước nghiệp vụ.
\end{itemize}

\subsubsection{FR-03 --- Clarification và giao dịch chờ xác nhận}

\textbf{Mục đích và kích hoạt.} Thu thập trường thiếu qua nhiều lượt và ngăn ghi dữ liệu suy đoán. FR-01/FR-02 tạo state khi thiếu trường, có nhiều lựa chọn hoặc đã đủ dữ liệu để xem trước.

\textbf{Đầu vào và xử lý.} State gắn với hồ sơ, gồm intent, trường chờ, dữ liệu đã thu, candidates và thời điểm tạo; pending chứa một giao dịch/batch cùng metadata. State và pending nằm trong KV store với TTL 5 phút. Câu trả lời được merge rồi validation lại; sửa preview dùng lấy--cập nhật--đặt có kiểm soát; xác nhận dùng \texttt{claim} nguyên tử để một \texttt{pending\_id} chỉ được tiêu thụ một lần.

\textbf{Lỗi, đầu ra và hậu điều kiện.} Pending hết hạn, sai ID hoặc đã claim không được ghi DB; sửa sai giữ preview cũ khi có thể; hủy xóa state; request cũ không ghi đè preview mới. Xác nhận hợp lệ xóa pending và chuyển dữ liệu cho transaction service.

\textbf{Tiêu chí chấp nhận.}
\begin{itemize}
  \item \texttt{TC-FR03-01}: câu trả lời bổ sung đúng trường thiếu và làm mới TTL.
  \item \texttt{TC-FR03-02}: hai xác nhận đồng thời chỉ một request claim và ghi dữ liệu.
  \item \texttt{TC-FR03-03}: ID cũ không tiêu thụ preview mới.
  \item \texttt{TC-FR03-04}: sửa sai hoặc hết TTL không đổi PostgreSQL.
\end{itemize}

\subsubsection{FR-04 --- Phân loại danh mục và học từ phản hồi}

\textbf{Mục đích và kích hoạt.} Chọn danh mục có thể giải thích, thích nghi với sửa sai cá nhân nhưng không tự học từ tín hiệu không chắc chắn. Matcher nhận mô tả, loại thu/chi, danh mục/alias và correction thuộc cùng hồ sơ.

\textbf{Đầu vào và xử lý.} Chuỗi được bỏ dấu, hạ chữ và chuẩn hóa khoảng trắng. Thứ tự là tên exact, alias exact rồi fuzzy theo mục 2.2.2; ngưỡng 0,90 cho input rất ngắn, 0,82 cho input còn lại và margin tối thiểu 0,08. Correction gần input có thể làm few-shot context; kết quả vẫn qua validation/xác nhận. Sửa category của giao dịch không-manual ghi cặp cũ--đúng best-effort sau commit.

\textbf{Lỗi, đầu ra và hậu điều kiện.} Không đạt ngưỡng, hai ứng viên quá gần, alias xung đột hoặc correction không đồng thuận phải về ``Khác''/yêu cầu chọn. Lỗi ghi feedback không biến giao dịch đã commit thành thất bại. Kết quả gồm category, confidence, kiểu match và lý do fallback.

\textbf{Tiêu chí chấp nhận.}
\begin{itemize}
  \item \texttt{TC-FR04-01}: exact và alias exact đứng trước fuzzy.
  \item \texttt{TC-FR04-02}: chênh lệch dưới 0,08 trả ambiguous/fallback.
  \item \texttt{TC-FR04-03}: input ngắn dưới 0,90 không được tự gán.
  \item \texttt{TC-FR04-04}: correction hợp lệ làm ví dụ; correction xung đột không ép phân loại.
\end{itemize}

\subsubsection{FR-05 --- Quản lý giao dịch, ví và danh mục}

\textbf{Mục đích và kích hoạt.} Duy trì sổ cái có thể sửa, soft delete, khôi phục và truy vấn mà không làm sai số dư. Người dùng CRUD giao dịch, ví và danh mục; bộ lọc hỗ trợ kỳ ngày, loại, danh mục, ví và phân trang.

\textbf{Đầu vào và xử lý.} Giao dịch thu--chi có amount dương, ngày không tương lai, category/wallet hợp lệ. Theo mục 2.1.1, income cộng và expense trừ số dư. Tạo, sửa amount/type/wallet, soft delete và restore đều áp phần chênh lệch trong cùng DB transaction; đổi category không đổi số dư; báo cáo loại bản ghi có \texttt{deleted\_at}.

\textbf{Lỗi, đầu ra và hậu điều kiện.} Dữ liệu sai miền hoặc tham chiếu ngoài hồ sơ bị từ chối; \textbf{số dư sau chi được phép âm}. Lỗi cache sau commit chỉ ghi log. API trả giao dịch hydrate và số dư mới; restore chỉ áp lại tác động đúng một lần.

\textbf{Tiêu chí chấp nhận.}
\begin{itemize}
  \item \texttt{TC-FR05-01}: income/expense làm số dư đổi lần lượt $+A/-A$.
  \item \texttt{TC-FR05-02}: sửa amount/type/wallet áp đúng chênh lệch và rollback toàn bộ khi lỗi.
  \item \texttt{TC-FR05-03}: expense lớn hơn số dư vẫn hợp lệ và ví có thể âm.
  \item \texttt{TC-FR05-04}: soft delete/restore đảo và áp số dư đúng một lần.
\end{itemize}

\subsubsection{FR-06 --- Chuyển ví và dòng tiền đặc biệt}

\textbf{Mục đích và kích hoạt.} Ghi transfer và dòng tiền đầu tư mà không tạo debit/credit dở dang. Amount phải dương, ngày không tương lai; transfer cần hai ví khác nhau, investment flow cần đầu ví phù hợp.

\textbf{Đầu vào và xử lý.} Service khóa ví theo thứ tự, trừ $A$ ở nguồn, cộng $A$ ở đích và ghi \texttt{wallet\_transfers} trong cùng \texttt{BEGIN/COMMIT}. Chuyển nội bộ giữ bất biến tổng thay đổi bằng 0 theo mục 2.1.1. Lãi/lỗ đầu tư khác 0 chỉ áp dụng cho ví investment/savings.

\textbf{Lỗi, đầu ra và hậu điều kiện.} Ví trùng/thiếu, amount/date sai hoặc ví ngoài hồ sơ làm rollback. \textbf{Không kiểm tra đủ số dư}; ví nguồn được phép âm. Trả lịch sử chuyển và hai số dư cùng phản ánh một lần chuyển hoặc cùng giữ nguyên.

\textbf{Tiêu chí chấp nhận.}
\begin{itemize}
  \item \texttt{TC-FR06-01}: $B'_s=B_s-A$, $B'_d=B_d+A$ và tổng số dư không đổi.
  \item \texttt{TC-FR06-02}: ví nguồn không đủ vẫn được phép âm.
  \item \texttt{TC-FR06-03}: hai ví trùng/ngoài hồ sơ bị từ chối trước commit.
  \item \texttt{TC-FR06-04}: lỗi giữa debit, credit và insert rollback cả ba hiệu ứng.
\end{itemize}

\subsubsection{FR-07 --- Phân tích dữ liệu tài chính}

\textbf{Mục đích và kích hoạt.} Biến lịch sử giao dịch thành dữ kiện định lượng về xu hướng, bất thường, nguy cơ cạn tiền, khoản chi lặp và liên hệ danh mục. Luồng chạy khi người dùng mở báo cáo/insight hoặc worker quét định kỳ.

\textbf{Đầu vào và xử lý.} Chỉ dùng giao dịch chưa xóa và zero-fill trục thời gian khi cần. Analytics Engine áp dụng mục 2.3.1--2.3.5: (i) trend chỉ công bố tăng với ít nhất 3 tháng, slope dương, tăng trung bình ít nhất 10\% và $R^2\geq0{,}5$; (ii) anomaly cần ít nhất 4 điểm, cờ phía chi lớn khi $z\geq2{,}5$ hoặc vượt $Q_3+1{,}5IQR$; (iii) runway dùng đúng 14 ngày lịch kể cả ngày chi 0, số dư không dương trả 0 ngày và burn rate 0 không dự báo ngày cạn; (iv) recurring dùng cửa sổ 90 ngày, amount tối đa 500.000~VND, sai số 15\%, cadence 20--40 ngày; (v) correlation zero-fill 12 tuần, cần dữ liệu ở ít nhất 4 tuần và chỉ công bố tương quan \textbf{dương} $r\geq0{,}6$, không dùng $|r|$.

\textbf{Lỗi, đầu ra và hậu điều kiện.} Thiếu dữ liệu trả \texttt{null}; nhánh biên xử lý mẫu số/burn rate bằng 0. Thành phần lỗi được ghi trong \texttt{degraded\_components}; correlation không được diễn giải là nhân quả. Facts JSON có thời điểm, giá trị, kỳ, phương pháp và cảnh báo; cache insight TTL 10 phút.

\textbf{Tiêu chí chấp nhận.}
\begin{itemize}
  \item \texttt{TC-FR07-01}: slope, forecast và $R^2$ khớp chuỗi tính tay.
  \item \texttt{TC-FR07-02}: z-score/IQR chỉ đánh dấu phía trên và xử lý phương sai/IQR bằng 0.
  \item \texttt{TC-FR07-03}: runway dùng đủ 14 ngày kể cả ngày chi 0; số dư âm trả 0 ngày.
  \item \texttt{TC-FR07-04}: recurring tuân số lần, amount và cadence; \texttt{TC-FR07-05}: chỉ cặp $r\geq0{,}6$ được trả, cặp âm không được công bố.
\end{itemize}

\subsubsection{FR-08 --- Insight có căn cứ và persona}

\textbf{Mục đích và kích hoạt.} Diễn giải FR-07 bằng ngôn ngữ dễ hiểu mà không giao phép tính cho LLM. Luồng nhận facts, kỳ dữ liệu, \texttt{degraded\_components} và persona; endpoint facts thô phải dùng độc lập với narrator.

\textbf{Đầu vào và xử lý.} Backend gửi facts có cấu trúc cho narrator, đồng thời tạo lời khuyên tổng quan bằng quy tắc xác định. Persona chỉ thay cách diễn đạt; response luôn kèm facts, provider và metadata để đối chiếu. Ranh giới ở mục 2.5.2, phép đo grounding ở mục 2.4.3.

\textbf{Lỗi, đầu ra và hậu điều kiện.} LLM không khả dụng thì dùng template fallback hoặc trả lỗi rõ ràng, không sinh số mock. \textbf{Checker bảo đảm mọi số trong narration thuộc facts mới là thiết kế đích}; chưa được tuyên bố numeric faithfulness đã đạt. Luồng đọc không ghi giao dịch, ngân sách hay mục tiêu.

\textbf{Tiêu chí chấp nhận.}
\begin{itemize}
  \item \texttt{TC-FR08-01}: đổi persona không đổi facts của cùng lần phân tích.
  \item \texttt{TC-FR08-02}: provider lỗi dùng fallback thật, giữ facts và không chèn dữ liệu mẫu.
  \item \texttt{TC-FR08-03} (thiết kế đích): checker phát hiện số/đơn vị không được facts hỗ trợ.
\end{itemize}

\subsubsection{FR-09 --- Ngân sách, đề xuất và dự báo}

\textbf{Mục đích và kích hoạt.} Theo dõi hạn mức, dự báo nguy cơ vượt và đề xuất dựa trên lịch sử thay vì để LLM tự chọn số. Người dùng CRUD ngân sách hoặc yêu cầu \texttt{category\_average}, \texttt{50-30-20}, \texttt{hybrid}.

\textbf{Đầu vào và xử lý.} Category phải là loại chi, limit dương, tháng/năm hợp lệ. Progress tính spent, remaining và tỷ lệ; forecast dùng $\text{spent}/\text{elapsed\_days}\times\text{days\_in\_month}$. Recommender theo mục 2.3.6 lấy trung bình trên toàn bộ tháng trong cửa sổ kể cả tháng 0, buffer 5\%; needs/wants có trần 50\%/30\% income và saving 20\%. Hybrid chỉ co tỷ trọng khi tổng vượt trần. Dưới 3 tháng sinh cảnh báo; apply cần \texttt{confirmed=true} và upsert nguyên tử.

\textbf{Lỗi, đầu ra và hậu điều kiện.} Thiếu income làm chiến lược tỷ lệ fallback về trung bình. Category/kỳ/limit sai, batch rỗng hoặc quá 50 phần tử hay thiếu xác nhận bị từ chối; một item sai rollback toàn batch. Chỉ endpoint apply đã xác nhận mới đổi \texttt{budgets}/\texttt{budget\_history}.

\textbf{Tiêu chí chấp nhận.}
\begin{itemize}
  \item \texttt{TC-FR09-01}: tháng không phát sinh vẫn nằm trong mẫu số trung bình.
  \item \texttt{TC-FR09-02}: forecast khớp công thức và không chia 0.
  \item \texttt{TC-FR09-03}: dưới 3 tháng có cảnh báo; thiếu income fallback đúng.
  \item \texttt{TC-FR09-04}: thiếu xác nhận bị chặn; item sai rollback batch.
\end{itemize}

\subsubsection{FR-10 --- Mục tiêu tài chính và mô phỏng what-if}

\textbf{Mục đích và kích hoạt.} Tính mức góp, thời gian hoàn thành và kịch bản cho mục tiêu tiết kiệm, mua sắm hoặc trả nợ bằng hàm xác định. Đầu vào có goal type, target dương, current/contribution không âm, ngày không quá khứ và lãi suất năm cho nợ.

\textbf{Đầu vào và xử lý.} Goal Planner áp dụng mục 2.3.7 để tính phần thiếu, số tháng, mức góp đến hạn, tiến độ và cảnh báo. Với nợ, lãi suất tháng bằng lãi suất năm chia 12; công thức niên kim có nhánh lãi suất 0 và mô phỏng tháng để tính lãi/remaining balance; khoản trả không vượt lãi tháng đầu bị cờ negative amortization. What-if chạy lại cùng hàm. Endpoint plan không ghi DB, cấp token ký theo payload trong 15 phút; tạo/cập nhật trường kế hoạch chỉ nhận token còn hạn và fingerprint khớp.

\textbf{Lỗi, đầu ra và hậu điều kiện.} Ngày/lãi suất/số tiền sai, token hết hạn hoặc payload bị đổi đều bị từ chối. Contribution 0 không có thời hạn; nợ không giảm trả cảnh báo. Preview trả plan, progress, cashflow, what-if và token nhưng không ghi dữ liệu; xác nhận hợp lệ mới lưu.

\textbf{Tiêu chí chấp nhận.}
\begin{itemize}
  \item \texttt{TC-FR10-01}: saving có/không deadline khớp công thức.
  \item \texttt{TC-FR10-02}: nợ lãi suất 0/dương khớp mô phỏng và phát hiện negative amortization.
  \item \texttt{TC-FR10-03}: preview không tạo row; token hết hạn/payload khác bị chặn.
  \item \texttt{TC-FR10-04}: what-if không thay goal đã lưu.
\end{itemize}

\subsubsection{FR-11 --- Khoản định kỳ và tác vụ chủ động}

\textbf{Mục đích và kích hoạt.} Quản lý lịch chi lặp, ghi thanh toán nhất quán và tạo nhắc/insight không trùng khi retry. Bill có tên, amount dương, frequency weekly/monthly/quarterly/yearly và due day hợp lệ; payment mang \texttt{period\_due\_date} vừa đọc; worker nhận job name và user scope.

\textbf{Đầu vào và xử lý.} Ngày đến hạn được tính theo chu kỳ và kẹp cuối tháng. Gợi ý recurring dùng mục 2.3.4. Payment khóa bill, kiểm tra kỳ, tạo expense, trừ ví, ghi payment và dời \texttt{next\_due\_date} trong một transaction; xóa plan vẫn giữ payment. Scheduler upsert lịch, thử tối đa 3 lần với exponential backoff; event key/fingerprint và unique index chống tạo message/export lặp.

\textbf{Lỗi, đầu ra và hậu điều kiện.} Kỳ thiếu/cũ, lịch/amount/reference sai làm rollback. Handler không hỗ trợ phải lỗi rõ; retry giữ cùng event key. Ví sau payment được phép âm. Trả bill, payment, giao dịch, kỳ kế tiếp và thống kê worker; một kỳ chỉ có một hiệu ứng.

\textbf{Tiêu chí chấp nhận.}
\begin{itemize}
  \item \texttt{TC-FR11-01}: bốn chu kỳ tạo đúng ngày đến hạn, kể cả tháng thiếu ngày 29--31.
  \item \texttt{TC-FR11-02}: payment cập nhật bốn hiệu ứng nguyên tử và rollback khi lỗi.
  \item \texttt{TC-FR11-03}: request kỳ cũ không thanh toán nhầm kỳ mới.
  \item \texttt{TC-FR11-04}: chạy lặp event key không tăng message/file lần hai.
\end{itemize}

\subsubsection{FR-12 --- Xuất dữ liệu, lịch sử tệp và dọn dẹp}

\textbf{Mục đích và kích hoạt.} Cho người dùng lấy dữ liệu theo bộ lọc, theo dõi tệp đã tạo và dọn tệp hết hạn. Chỉ giao dịch chưa soft delete, thuộc hồ sơ và thỏa format/khoảng ngày/filter được xuất.

\textbf{Đầu vào và xử lý.} CSV UTF-8 có BOM, escape ký tự đặc biệt và lưu history với TTL mặc định 7 ngày. Luồng API mang tên \texttt{pdf} hiện tạo \texttt{.html}, trả \texttt{text/html}; đây là HTML có thể in, \textbf{không phải PDF nhị phân}. Worker xóa file quá hạn và cập nhật trạng thái history. Endpoint backup \texttt{.pfbak} tồn tại nhưng chưa được kiểm chứng end-to-end nên không phải bằng chứng phục hồi production.

\textbf{Lỗi, đầu ra và hậu điều kiện.} Không có dữ liệu phù hợp thì không tạo file rỗng như thành công; format/filter sai, file hết hạn hoặc ID ngoài hồ sơ trả lỗi; không trả URL giả khi ghi file lỗi. Cleanup làm file không tải được nhưng giữ metadata theo chính sách.

\textbf{Tiêu chí chấp nhận.}
\begin{itemize}
  \item \texttt{TC-FR12-01}: CSV lọc đúng, loại soft-delete và escape ký tự đặc biệt.
  \item \texttt{TC-FR12-02}: endpoint nhãn PDF trả \texttt{.html}/\texttt{text/html}, không được gọi là PDF thật.
  \item \texttt{TC-FR12-03}: file quá TTL bị dọn và history báo không khả dụng.
  \item \texttt{TC-FR12-04}: chỉ công bố backup/restore sau test toàn vẹn, checksum, rollback và đối soát bảng.
\end{itemize}

\subsection{Yêu cầu phi chức năng}

Các ngưỡng dưới đây là tiêu chí nghiệm thu hoặc thiết kế đích, không mặc nhiên là kết quả hiện tại. Kết quả thực đo và phép đo còn thiếu được tách tại mục 3.3.2.

\begin{longtable}{@{}L{1.2cm}L{2.8cm}L{3.5cm}L{5.2cm}@{}}
  \caption{Yêu cầu phi chức năng và cách đo}\label{tab:nfr}\\
  \toprule
  \textbf{Mã} & \textbf{Yêu cầu} & \textbf{Mục tiêu} & \textbf{Cách đo} \\
  \midrule
  \endfirsthead
  \toprule
  \textbf{Mã} & \textbf{Yêu cầu} & \textbf{Mục tiêu} & \textbf{Cách đo} \\
  \midrule
  \endhead
  NFR-01 & Tính đúng và nguyên tử dữ liệu & Không partial write ở create, update, delete, restore, transfer, recurring payment và batch apply; số dư âm hợp lệ không bị tính là lỗi. & Fault injection từng bước, đối chiếu row/số dư trước--sau và kiểm tra rollback. \\
  NFR-02 & Độ chính xác trích xuất & Khóa ngưỡng riêng cho amount, type, date, category, wallet trên tập gán nhãn độc lập; không suy diễn từ 31 ca hard-coded. & Precision, recall, F1 từng field, exact match và phân tích lỗi theo kiểu câu. \\
  NFR-03 & Tính trung thực số liệu & \textbf{Thiết kế đích:} narration không thêm/đổi số hoặc đơn vị ngoài facts; chưa tuyên bố đã đạt. & Checker facts--response, numeric faithfulness và unsupported-number rate trên tập khóa phiên bản. \\
  NFR-04 & Hiệu năng & Text p95 mục tiêu $\leq3$ giây; media p95 $\leq8$ giây trong môi trường công bố. & Ít nhất 30 lượt/luồng; ghi phần cứng, provider/model, cache hit/miss, p50/p95 và error rate. \\
  NFR-05 & Khả năng phục hồi & Provider lỗi không tạo dữ liệu giả; Redis lỗi có fallback phát triển; lỗi hậu commit không gây retry trùng. & Test timeout/unavailable cho LLM/OCR/STT/Redis, kiểm tra response, pending và DB. \\
  NFR-06 & Riêng tư và phạm vi bảo mật & Không ghi credential; xóa media tạm; traits chỉ dùng khi consent. & Kiểm tra log/config/file tạm và consent; ghi rõ \texttt{default\_user} chưa cô lập multi-user production. \\
  NFR-07 & Khả năng kiểm thử & Giải thuật thuần không phụ thuộc DB/LLM; ca thường và biên có expected tính tay. & Unit test trend, anomaly, runway, correlation, budget, goal, matching và báo cáo coverage/truy vết. \\
  NFR-08 & Nhất quán tác vụ nền & Retry cùng event/job không tạo message/export hoặc thanh toán lặp. & Chạy cùng fingerprint nhiều lần, giả lập lỗi giữa bước và so bản ghi; Redis worker live đo riêng. \\
  NFR-09 & Khả năng truy vết & Insight trả facts, kỳ dữ liệu, phương pháp, provider và thành phần suy giảm. & Contract test và đối chiếu FR $\rightarrow$ công thức $\rightarrow$ test $\rightarrow$ artifact. \\
  NFR-10 & Bảo trì và thay thế & Module route--service--model, adapter provider và ngưỡng cấu hình không đổi contract nghiệp vụ. & Review phụ thuộc, test thay provider/fallback và kiểm tra không đặt công thức trong route/narrator. \\
  \bottomrule
\end{longtable}

Các giới hạn nghiệm thu hiện tại gồm: chưa có authentication/authorization production; grounding checker mới là thiết kế đích; Redis worker chưa được smoke test live đầy đủ; luồng ``PDF'' mới trả HTML; backup/restore chưa đủ bằng chứng để tuyên bố hoàn chỉnh.

\section{Thiết kế phần mềm}

\subsection{Kiến trúc ứng dụng}

PERFIN dùng modular monolith thay vì microservice. API Express chứa các module nghiệp vụ độc lập ở mức mã nguồn nhưng cùng tiến trình triển khai. Worker là tiến trình riêng cho job nền. Cấu trúc này phù hợp quy mô niên luận, giảm chi phí vận hành nhưng vẫn giữ ranh giới module.

\widereportfigure{02-runtime-architecture.pdf}{Kiến trúc vận hành của PERFIN. Lõi xác định tạo và lưu facts; AI Orchestrator chỉ điều phối ngữ nghĩa.}{fig:runtime-architecture}

Hình~\ref{fig:runtime-architecture} thể hiện đường đi chuẩn: mobile gọi route, route chuyển cho service, service dùng model truy cập PostgreSQL. Analytics Engine, Budget Engine và Goal Planner tính kết quả xác định. Redis phục vụ cache/state/queue; worker xử lý tác vụ định kỳ; provider AI nằm ngoài biên tin cậy dữ liệu.

\begin{table}[H]
  \centering
  \caption{Thành phần kiến trúc và trách nhiệm}
  \label{tab:components}
  \begin{tabularx}{\textwidth}{@{}L{3.4cm}Y Y@{}}
    \toprule
    \textbf{Thành phần} & \textbf{Trách nhiệm chính} & \textbf{Không chịu trách nhiệm} \\
    \midrule
    Mobile App & Thu input, preview, hiển thị và xác nhận. & Tính insight hoặc giữ bí mật provider. \\
    Express Routes & HTTP contract, validation đầu vào, mapping response. & Chứa công thức phân tích. \\
    Core Services & Luật nghiệp vụ, transaction và xác nhận. & Sinh câu trả lời tùy ý. \\
    AI Orchestrator & Tool routing, extraction và fallback. & Truy cập DB không qua tool. \\
    Analytics Engine & Tính facts thống kê. & Viết lời khuyên cá nhân hóa. \\
    Persona/Narrator & Diễn đạt facts theo giọng điệu. & Thay số hoặc quyết định. \\
    PostgreSQL Models & Truy vấn, constraint và transaction. & Lưu state hội thoại tạm. \\
    Redis/KV & TTL state, cache và rate count. & Làm nguồn dữ liệu tài chính chuẩn. \\
    BullMQ Worker & Job định kỳ, retry và dedup. & Xử lý request tương tác trực tiếp. \\
    \bottomrule
  \end{tabularx}
\end{table}

\widereportfigure{03-deployment.pdf}{Sơ đồ triển khai nguyên mẫu. Frontend, API, worker, PostgreSQL và Redis thuộc môi trường demo; provider AI được gọi qua mạng.}{fig:deployment}

Hình~\ref{fig:deployment} là \textbf{thiết kế triển khai nguyên mẫu}, không mô tả cụm production hay cam kết sẵn sàng cao. Docker Compose hiện chỉ hỗ trợ Redis; API, worker và PostgreSQL cần được khởi chạy/cấu hình riêng. Readiness phải tách khỏi liveness để backend không được coi là sẵn sàng khi bootstrap DB thất bại.

\subsection{Thiết kế dữ liệu}

\subsubsection{Mô hình miền}

\widereportfigure{04-domain-class.pdf}{Mô hình miền theo aggregate nghiệp vụ. Các cụm dữ liệu tách khỏi service phân tích và lớp ngôn ngữ.}{fig:domain-class}

Hình~\ref{fig:domain-class} tách bốn aggregate sổ cái, lập kế hoạch, khoản định kỳ và hội thoại--phản hồi. Hộp User được lặp làm neo bố cục nhưng cùng chỉ một thực thể; vì vậy quan hệ sở hữu và luồng đọc facts đều là đường thẳng, không giao cắt. Analytics/Budget/Goal Services chỉ đọc aggregate để tính facts hoặc kế hoạch và không sở hữu entity tài chính.

\subsubsection{Mô hình vật lý}

\widereportfigure{05-physical-erd.pdf}{Sơ đồ quan hệ thực thể vật lý được đối chiếu với chuỗi migration runtime.}{fig:physical-erd}

ERD ở Hình~\ref{fig:physical-erd} gồm 18 bảng vật lý: \code{users}, \code{ai_personalities}, \code{user_traits}, \code{categories}, \code{wallets}, \code{transactions}, \code{budgets}, \code{budget_history}, \code{chat_messages}, \code{investment_pnl}, \code{wallet_transfers}, \code{export_history}, \code{backup_config}, ba bảng recurring, \code{financial_goals} và \code{ai_feedback_logs}. Ba compatibility view không được đếm như bảng. Một bảng \code{users} được lặp thành bốn neo tham chiếu; FK chéo mô-đun được ghi trong hộp bảng con thay vì kéo đường dài. Cách trình bày này giữ PK/FK/UQ/CHECK nhưng loại đường cong và phần lớn giao cắt. Khóa \code{users.user_key} là cầu nối với \code{default_user}, không phải bằng chứng đã có authentication.

\begin{table}[H]
  \centering
  \caption{Nhóm thực thể dữ liệu chính}
  \label{tab:data-entities}
  \begin{tabularx}{\textwidth}{@{}L{3.8cm}Y Y@{}}
    \toprule
    \textbf{Nhóm} & \textbf{Bảng/thực thể} & \textbf{Ý nghĩa} \\
    \midrule
    Người dùng và cá nhân hóa & \code{users}, \code{ai_personalities}, \code{user_traits} & Hồ sơ demo, persona, consent và traits. \\
    Sổ cái cá nhân & \code{wallets}, \code{categories}, \code{transactions} & Nguồn thu--chi cốt lõi. \\
    Ngân sách & \code{budgets}, \code{budget_history} & Hạn mức và lịch sử thay đổi. \\
    Chi định kỳ & Ba bảng recurring & Lịch, thanh toán và gợi ý bị bỏ qua. \\
    Phản hồi AI & \code{ai_feedback_logs} & Kết quả ban đầu và sửa của người dùng. \\
    Mục tiêu & \code{financial_goals} & Tiết kiệm, mua sắm, trả nợ và tiến độ. \\
    Dòng tiền đặc biệt & \code{wallet_transfers}, \code{investment_pnl} & Chuyển ví và lãi/lỗ đầu tư. \\
    Vận hành & \code{chat_messages}, \code{export_history}, \code{backup_config} & Chat, export và cấu hình sao lưu. \\
    \bottomrule
  \end{tabularx}
\end{table}

\subsubsection{Quy tắc dữ liệu quan trọng}

\begin{itemize}
  \item Tiền dùng \code{DECIMAL(15,2)}; service phải chuyển đổi rõ kiểu khi tính trong JavaScript.
  \item Soft delete giữ khả năng phục hồi; truy vấn báo cáo mặc định loại bản ghi đã xóa.
  \item Xóa danh mục có giao dịch phải bị chặn hoặc re-tag có kiểm soát, không cascade làm mất lịch sử.
  \item Pending/clarification nằm trong Redis với schema và TTL riêng, không thuộc ERD bền vững.
  \item Mọi truy vấn theo ID trong thiết kế đích phải kèm \code{user_id/user_key} để ngăn truy cập chéo khi bổ sung xác thực.
\end{itemize}

\subsection{Thiết kế chi tiết}

\subsubsection{Ranh giới LLM trong kiến trúc}

\reportfigure[0.88\linewidth]{06-llm-boundary.pdf}{Ranh giới trách nhiệm của LLM: dữ liệu và facts được tạo ở lõi xác định; LLM chỉ chuyển đổi ngữ nghĩa và diễn giải.}{fig:llm-boundary}

Hình~\ref{fig:llm-boundary} là sơ đồ quan trọng nhất để giải thích tác dụng LLM. Đầu vào tự nhiên được biến thành typed command; service kiểm tra và tạo preview. Với câu hỏi phân tích, SQL và giải thuật tạo facts trước khi narrator nhận dữ liệu. LLM không có kết nối trực tiếp PostgreSQL và không tự thực thi tool. Lợi ích của LLM được đánh giá bằng độ chính xác thực thể, tỷ lệ hỏi lại, số thao tác tiết kiệm và chất lượng diễn giải, không bằng việc thay thế công thức.

\subsubsection{Máy trạng thái hội thoại}

\reportfigure[0.86\linewidth]{07-conversation-state.pdf}{Máy trạng thái hội thoại và giao dịch chờ xác nhận.}{fig:conversation-state}

Các trạng thái chính ở Hình~\ref{fig:conversation-state} gồm idle, parse, collecting, preview, confirmed, cancelled và expired; lựa chọn ứng viên được giữ trong collecting. Redis giữ intent, trường đang chờ, dữ liệu đã thu và candidates trong một TTL giới hạn. Nhánh ghi dữ liệu bắt buộc đi qua preview rồi confirmed. Ba trạng thái cuối kết thúc phiên hiện tại; yêu cầu mới tạo phiên idle khác. Thiết kế đích dùng thao tác claim nguyên tử khi xác nhận để hai request đồng thời không cùng tạo giao dịch.

\subsubsection{Nhập giao dịch bằng văn bản}

\widereportfigure{08-text-sequence.pdf}{Sơ đồ tuần tự nhập giao dịch bằng văn bản. LLM nằm ngoài transaction ghi dữ liệu.}{fig:text-sequence}

Trình tự Hình~\ref{fig:text-sequence} gồm nhận text, lấy categories/wallets và corrections, gọi LLM tool hoặc local parser, chuẩn hóa, validation, hỏi lại nếu thiếu, tạo preview, sửa/xác nhận và ghi DB trong transaction. Kiểm tra ngân sách diễn ra sau khi số dư và giao dịch đã nhất quán. Nếu commit lỗi, không được trả thông báo thành công.

\subsubsection{Đầu vào đa phương thức}

\reportfigure[0.84\linewidth]{09-multimodal-flow.pdf}{Luồng xử lý đầu vào đa phương thức. Voice và ảnh hội tụ vào pipeline text chung sau bước xác nhận.}{fig:multimodal-flow}

Voice đi qua STT và bước xác nhận transcript. Ảnh đi qua preprocessing/OCR rồi cho phép chọn tổng hóa đơn hoặc từng mặt hàng. Hai nhánh hội tụ tại pipeline text. Lỗi provider phải được biểu diễn là lỗi hoặc yêu cầu nhập tay; không dùng raw text giả. Báo cáo hiện chưa có ground truth media nên Hình~\ref{fig:multimodal-flow} mô tả \textbf{thiết kế đích}, không phải bằng chứng accuracy OCR/STT.

\subsubsection{Phân loại và vòng phản hồi}

\reportfigure[0.84\linewidth]{10-feedback-flow.pdf}{Luồng phản hồi và cá nhân hóa phân loại bằng correction retrieval, không phải fine-tuning.}{fig:feedback-flow}

Khi người dùng sửa category, hệ thống lưu cặp kết quả AI--kết quả đúng. Lần sau, correction exact/fuzzy được xét trước, sau đó mới đến matcher/LLM. Giao dịch lặp trong ``Khác'' có thể được gom cụm để đề xuất category; tạo category và re-tag luôn là kế hoạch chờ xác nhận. Lịch sử mâu thuẫn phải hạ confidence thay vì học theo correction gần nhất.

\subsubsection{Sinh insight có căn cứ}

\widereportfigure{11-insight-sequence.pdf}{Sơ đồ tuần tự sinh insight có căn cứ. Response trả cả thông điệp và facts để truy vết.}{fig:insight-sequence}

Trong Hình~\ref{fig:insight-sequence}, model SQL truy xuất dữ liệu; Analytics Engine tính facts và metadata; guard kiểm tra đơn vị cùng số quan sát. Chỉ sau đó narrator mới diễn giải theo persona. Khi LLM lỗi, narrator template dùng chính facts. Thiết kế đích bổ sung numeric checker trước khi trả response để phát hiện số không được hỗ trợ.

\subsubsection{Lập kế hoạch mục tiêu}

\reportfigure[0.86\linewidth]{12-goal-flow.pdf}{Luồng lập kế hoạch mục tiêu và mô phỏng what-if.}{fig:goal-flow}

LLM trong Hình~\ref{fig:goal-flow} chỉ trích xuất mục tiêu, số tiền, hạn chót, mức góp và lãi suất. Goal Planner kiểm tra miền dữ liệu, tính surplus, thời hạn, mức góp và cảnh báo. What-if tạo kịch bản mới mà không sửa kế hoạch gốc. Chỉ xác nhận mới lưu \code{financial_goals}.

\subsubsection{Tác vụ chủ động}

\widereportfigure{13-worker-sequence.pdf}{Sơ đồ tuần tự tác vụ chủ động và cơ chế chống hiệu ứng lặp khi retry.}{fig:worker-sequence}

Scheduler tạo job có định danh; worker lấy đúng user scope, gọi handler, tính facts và ghi internal message hoặc export. Retry dùng cùng fingerprint; unique event key và thao tác idempotent loại thông báo trùng. Tính năng hiện tại tạo thông báo trong \code{chat_messages}, không phải push notification. Auto-backup theo lịch chưa có worker tương ứng nên không được mô tả là đã hoạt động.

\subsubsection{Kiểm soát thao tác do LLM khởi tạo}

\begin{longtable}{@{}L{3.0cm}L{5.0cm}L{1.5cm}L{3.2cm}@{}}
  \caption{Điều kiện kiểm soát thao tác do LLM khởi tạo}\label{tab:llm-actions}\\
  \toprule
  \textbf{Thao tác} & \textbf{Validation bắt buộc} & \textbf{Xác nhận} & \textbf{Nguyên tử/idempotent} \\
  \midrule
  \endfirsthead
  \toprule
  \textbf{Thao tác} & \textbf{Validation bắt buộc} & \textbf{Xác nhận} & \textbf{Nguyên tử/idempotent} \\
  \midrule
  \endhead
  Tạo một/nhiều giao dịch & amount $>0$, type/category/wallet hợp lệ & Có & Batch và số dư cùng transaction. \\
  Chuyển ví & Hai ví khác nhau, amount/date hợp lệ; số dư âm được phép & Có & Debit, credit, history cùng transaction. \\
  Tạo recurring bill & Tên, số tiền, chu kỳ và ngày hợp lệ & Có & Không tạo trùng cùng kế hoạch. \\
  Áp dụng ngân sách & Lịch sử và tổng hạn mức hợp lệ & Có & Upsert theo user/category/kỳ. \\
  Tạo mục tiêu & Type, target, date, interest phù hợp & Có & Lưu sau preview. \\
  Tạo category và re-tag & Tên không generic, ID thuộc user & Có & Tạo/reuse, retag và feedback cùng kế hoạch. \\
  Truy vấn/insight & Khoảng dữ liệu hợp lệ & Không & Facts truy vết được. \\
  Export & Định dạng và khoảng ngày hợp lệ & Có từ chat & Cleanup có fingerprint. \\
  \bottomrule
\end{longtable}

\subsection{Bảo mật, riêng tư và đạo đức}

Nguyên mẫu chưa có authentication production; vì vậy không được tuyên bố đã cô lập nhiều người dùng. Các route dùng \code{default_user} và một số thao tác theo ID chưa kèm user scope. Trước triển khai thật cần bổ sung xác thực, authorization ở mọi truy vấn và test truy cập chéo.

Credential AI, service account và chuỗi kết nối không được ghi vào Git hoặc chat log. Ảnh/audio cần chính sách xóa sau xử lý. \code{user_traits} chỉ được dùng khi consent bật. Persona không được gây áp lực, miệt thị hoặc biến gợi ý thành tư vấn đầu tư chắc chắn. Backup không được dùng khóa AES đóng gói sẵn trong client; khóa phải đến từ cơ chế quản lý bí mật bên ngoài ứng dụng.

\section{Kiểm thử}

\subsection{Kế hoạch kiểm thử}

\subsubsection{Ma trận kiểm thử}

\begin{longtable}{@{}L{2.1cm}L{3.5cm}L{3.5cm}L{3.5cm}@{}}
  \caption{Ma trận kiểm thử}\label{tab:test-matrix}\\
  \toprule
  \textbf{Cấp} & \textbf{Đối tượng} & \textbf{Kỹ thuật} & \textbf{Bằng chứng} \\
  \midrule
  \endfirsthead
  \toprule
  \textbf{Cấp} & \textbf{Đối tượng} & \textbf{Kỹ thuật} & \textbf{Bằng chứng} \\
  \midrule
  \endhead
  Unit & Analytics, matching, budget, goal, validation & Dữ liệu thường, biên, invariant & \code{node:test}, expected tính tay. \\
  Integration & Route--service--model, Redis, PostgreSQL & DB test, rollback và TTL & Log chạy và snapshot DB. \\
  Contract AI & Tool schema và parser & Tập câu gán nhãn, provider/fallback & JSON từng case. \\
  Media & OCR/STT adapter & Ảnh/audio thật có ground truth & CER/WER và field accuracy. \\
  Grounding & Narrator/persona & So số response với facts & Numeric faithfulness. \\
  Worker & Retry, chống trùng và lịch chạy & Chạy lặp cùng event, giả lập lỗi & Số hiệu ứng trước--sau. \\
  System & Text/voice/image tới DB & End-to-end trên thiết bị & Video/log và DB assertion. \\
  UAT & Người dùng mục tiêu & Nhiệm vụ có thời gian/số bước & Tỷ lệ hoàn thành, SUS. \\
  \bottomrule
\end{longtable}

\subsubsection{Thiết kế dữ liệu kiểm thử}

Bộ dữ liệu văn bản phải tách tập phát triển và đánh giá, bao phủ các đơn vị \texttt{k}, \texttt{nghìn}, \texttt{ngàn}, \texttt{tr}, \texttt{triệu}, số bằng chữ, phép nhân, ngày tương đối, thu--chi--chuyển--đầu tư, nhiều giao dịch, typo, Việt--Anh, trường hợp thiếu/mơ hồ và correction xung đột. Tập media phải lưu ảnh/audio gốc cùng transcript hoặc bảng hóa đơn chuẩn. Dữ liệu analytics nên là dữ liệu tổng hợp/ẩn danh có slope, outlier, cadence và correlation biết trước để đối chiếu công thức.

\subsubsection{Bộ chỉ số}

\begin{table}[H]
  \centering
  \caption{Bộ chỉ số đánh giá}
  \label{tab:metrics}
  \begin{tabularx}{\textwidth}{@{}L{3.2cm}L{5.2cm}Y@{}}
    \toprule
    \textbf{Nhóm} & \textbf{Chỉ số} & \textbf{Cách diễn giải} \\
    \midrule
    Extraction & Precision, recall, F1 theo field, exact match & So amount/type/date/category/wallet với nhãn. \\
    Classification & Accuracy, macro-F1, confusion matrix & Không để danh mục nhiều mẫu che danh mục ít mẫu. \\
    Clarification & Completion rate, số lượt hỏi & Chỉ tính case ban đầu thiếu/mơ hồ. \\
    OCR/STT & CER/WER và transaction-field accuracy & Tách chất lượng raw và tác động nghiệp vụ. \\
    Grounding & Numeric faithfulness, unsupported-number rate & Mọi số phải map về facts. \\
    Analytics & Sai số so expected & Test công thức và điều kiện biên. \\
    Hiệu năng & p50, p95, error rate & Tách cache hit/miss và provider. \\
    Worker & Duplicate-effect rate, retry success & Lặp cùng fingerprint. \\
    Trải nghiệm & Hoàn thành, thời gian, số thao tác & So chat với form cùng nhiệm vụ. \\
    \bottomrule
  \end{tabularx}
\end{table}

\subsection{Kết quả và phân tích}

\subsubsection{Kết quả đã đo}

Ngày 16/07/2026, các phép kiểm tra từ unit đến runtime được chạy trong workspace. Kết quả được trình bày đúng phạm vi bằng chứng, không mở rộng thành kết luận về độ chính xác mô hình hoặc mức sẵn sàng production.

\begin{longtable}{@{}L{3.3cm}L{4.2cm}L{1.8cm}L{3.4cm}@{}}
  \caption{Kết quả đã đo và phép đo còn thiếu}\label{tab:measured-results}\\
  \toprule
  \textbf{Hạng mục} & \textbf{Kết quả} & \textbf{Phân loại} & \textbf{Phạm vi kết luận} \\
  \midrule
  \endfirsthead
  \toprule
  \textbf{Hạng mục} & \textbf{Kết quả} & \textbf{Phân loại} & \textbf{Phạm vi kết luận} \\
  \midrule
  \endhead
  Backend \code{npm test} sau sửa & 100/100 test đạt, 0 thất bại & \statusmeasured & Unit/service; gồm transaction, concurrency, recurring, health, importer và time-series. \\
  Local parser quality gate & 31/31 strict pass; 0 partial; 0 fail & \statusmeasured & Cục bộ, không dùng Gemini; strict exit code, 31 ca hard-coded. \\
  Baseline trước sửa & 13/13 tệp test; parser 27/31 strict & \statusmeasured & Giữ để phân tích lỗi và chứng minh tác động bản sửa. \\
  Full API/DB/media smoke & 23/23 kiểm tra đạt & \statusmeasured & PostgreSQL live, preview, sửa, hủy, OCR 2 ảnh và STT 1 M4A; không gồm Redis worker. \\
  Mobile-web UI smoke & 4/4 cổng đạt tại 390$\times$844 & \statusmeasured & Ba màn hình không overflow; ảnh tải lên hiển thị thật; Chromium. \\
  Expo web/Android export & Đóng gói 653/960 module & \statusmeasured & Chứng minh frontend đóng gói được; không phải hiệu năng runtime. \\
  Import CSV tài chính & 5.265 dòng, 0 reject, provenance đầy đủ & \statusmeasured & Đã chạy hai lần trên clone và đối soát PostgreSQL live sau backup. \\
  LLM trên tập gán nhãn & Chưa công bố & \statusmissing & Cần khóa dataset, model, prompt và cấu hình provider. \\
  OCR/STT accuracy & Chưa công bố & \statusmissing & Cần ảnh/audio cùng ground truth. \\
  Redis worker đang chạy & Chưa xác nhận trong đợt đo & \statusmissing & WSL không có Redis/Docker; API fallback và readiness degraded đã được kiểm tra. \\
  Numeric grounding, p50/p95, UAT & Chưa công bố & \statusmissing & Cần checker, môi trường và kịch bản khóa phiên bản. \\
  \bottomrule
\end{longtable}

Baseline parser có hai lỗi amount: ``1 triệu 5'' và phép nhân ``3 cái áo, mỗi cái 200k''; hai ca partial liên quan ``quần jeans'' và ``đi ăn''. Sau sửa, cả bốn ca có regression test và harness đạt 31/31 strict, đồng thời trả mã thoát khác 0 khi bất kỳ ca nào không đạt. Tỷ lệ 100\% chỉ áp dụng cho tập cố định này, chưa phải benchmark Gemini hoặc tập gán nhãn độc lập.

Full smoke chứng minh luồng HTTP--PostgreSQL, Gemini cấu trúc đầu ra, OCR và STT có thể hoàn tất trên các fixture đã chọn. Kết quả này chỉ là smoke test khả dụng: không có ground truth đủ để suy ra accuracy, CER/WER hoặc numeric faithfulness; một lần đo thời gian cũng không phải p50/p95. Redis worker chưa được chạy live trong môi trường ghi nhận.

\subsubsection{Trạng thái chức năng trọng yếu}

\begin{longtable}{@{}L{3.2cm}L{2.5cm}L{3.5cm}L{3.5cm}@{}}
  \caption{Trạng thái và hành vi đích của chức năng trọng yếu}\label{tab:feature-status}\\
  \toprule
  \textbf{Chức năng} & \textbf{Hiện trạng} & \textbf{Khoảng trống chính} & \textbf{Điều kiện hoàn tất} \\
  \midrule
  \endfirsthead
  \toprule
  \textbf{Chức năng} & \textbf{Hiện trạng} & \textbf{Khoảng trống chính} & \textbf{Điều kiện hoàn tất} \\
  \midrule
  \endhead
  Giao dịch và analytics & \statusmeasured & Transaction, post-commit, zero-fill và full DB smoke đã đạt; 5.265 dòng live đã đối soát. & Dataset gán nhãn và fault-injection DB rộng hơn. \\
  Chat preview/confirm & \statusmeasured & Claim một lần, stale ID, edit/cancel race, history và HTTP--DB smoke đã đạt. & Đo TTL/claim với Redis thật và tải đồng thời cao hơn. \\
  Local parser & \statusmeasured & Đạt 31/31 strict trên quality gate cố định. & Tách tập độc lập, tăng biến thể và báo macro-F1. \\
  LLM tool routing & \statusimplemented & Provider selection/media fallback đã sửa; chưa benchmark. & Contract test trên provider thật và tập gán nhãn. \\
  OCR/STT & \statusmeasured & Smoke provider thật đạt với 2 ảnh và 1 M4A; chưa có ground truth. & Adapter failure test, CER/WER và field accuracy trên tập ẩn danh. \\
  Ví/chuyển ví & Hiện thực một phần & Thiếu route/UI tạo ví trên fresh install; chưa có integration test transfer. & CRUD tối thiểu và kiểm thử toàn vẹn số dư. \\
  Recurring và worker & \statusmeasured & Payment đã khóa hàng, nguyên tử, có kỳ dự kiến và chat preview; notification vẫn là chat nội bộ. & Live DB test và kênh nhắc ngoài ứng dụng nếu mở rộng. \\
  Export/backup & Chưa đúng tên gọi & Mẫu số/escaping đã sửa; ``PDF'' vẫn là HTML, backup còn thiếu. & PDF thật hoặc đổi nhãn; backup/restore an toàn. \\
  Auth/multi-user & Ngoài phạm vi hiện tại & \code{default_user}, ID chưa luôn scoped. & Không demo như đã có; thêm trước production. \\
  \bottomrule
\end{longtable}

\subsubsection{Thí nghiệm cần thực hiện}

Ba thí nghiệm có giá trị cao nhất cho bản hoàn thiện là:

\begin{enumerate}
  \item \textbf{Ablation LLM so với parser cục bộ:} cùng tập câu, so field F1, clarification rate, latency và chi phí.
  \item \textbf{Feedback before/after:} phát lại nhóm câu sau corrections, đo category accuracy và kiểm tra không suy giảm nhóm khác.
  \item \textbf{Grounded narration:} cấp cùng facts cho nhiều persona, kiểm tra số liệu giữ nguyên trong khi giọng điệu thay đổi.
\end{enumerate}

Kết quả chỉ được chuyển vào tóm tắt và kết luận sau khi có log, cấu hình và dữ liệu khóa phiên bản.
